% Original Prism narrative-review starter for Drug Discovery Today.
% Elsevier's distribution-provided elsarticle class remains unchanged.
% This is an editable submission draft, not a claim of final journal compliance.
\documentclass[review,12pt]{elsarticle}
\usepackage{amsmath,amssymb,graphicx,array,booktabs,tabularx,siunitx}
\usepackage[hidelinks]{hyperref}
\sisetup{uncertainty-mode=separate}
\setcitestyle{super,sort&compress}
\journal{Drug Discovery Today}

\begin{document}
\begin{frontmatter}
\title{Review title}
\author[inst1]{First Author}
\author[inst1,inst2]{Second Author\texorpdfstring{\corref{corresponding}}{}}
\cortext[corresponding]{Corresponding author.}
\ead{second.author@example.org}
\address[inst1]{Department, First Institution, City, Country}
\address[inst2]{Department, Second Institution, City, Country}

\begin{abstract}
Replace this paragraph with a concise overview of the topic, scope, relevant
published evidence and implications. Adapt the abstract to the journal's
current instructions. This starter contains synthetic layout examples only;
it reports no scientific result.
\end{abstract}
\begin{keyword}
keyword one \sep keyword two \sep keyword three
\end{keyword}
\end{frontmatter}

\section{Introduction}\label{sec:introduction}
Introduce the question, explain why it matters and define the scope of the
review. The citation here tests bibliography mechanics only; retain it only
if relevant to the manuscript\cite{wilkinson2016fair}. Replace all examples,
placeholder information and invented numerical values before submission.

\section{Background and concepts}\label{sec:background}
Explain the concepts readers need to understand the review. Organize the
text around the topic, with references to the actual literature supplied by
the authors. Define key terminology and avoid treating an illustration as
research evidence.

\section{Current approaches and evidence}\label{sec:evidence}
\subsection{A thematic comparison}
Compare the approaches relevant to the question and explain the evidence
supporting each interpretation. Figure~\ref{fig:overview} illustrates a
reusable two-panel figure; replace it with graphics created for the review.

\begin{figure}[htbp]
\centering
\begin{minipage}[t]{0.47\linewidth}
\centering\textbf{a}\par\smallskip
\fbox{\parbox[c][32mm][c]{0.90\linewidth}{\centering
Overview graphic\\[5pt]Replace with an original PDF or PNG}}
\end{minipage}\hfill
\begin{minipage}[t]{0.47\linewidth}
\centering\textbf{b}\par\smallskip
\fbox{\parbox[c][32mm][c]{0.90\linewidth}{\centering
Comparison graphic\\[5pt]Keep labels readable at final size}}
\end{minipage}
\caption{\textbf{Two-panel layout example.} (a) Overview slot. (b) Comparison
slot. Replace the caption with the actual figure's description, including
units, symbols and uncertainty as applicable. These boxes represent no data.}
\label{fig:overview}
\end{figure}

\subsection{Comparison across studies}
Table~\ref{tab:comparison} shows a way to align invented numerical values.
Replace the cells with the authors' verified comparisons and cite their
sources. When a comparison is not numerical, use the same layout with text.
\begin{table}[htbp]
\centering
\caption{Synthetic comparison used solely to demonstrate table layout.}
\label{tab:comparison}
\small
\begin{tabularx}{\linewidth}{@{}>{\raggedright\arraybackslash}X
 S[table-format=3.0] S[table-format=1.3(3)] S[table-format=1.3]@{}}
\toprule
Example & {Count} & {Estimate $\pm$ SD} & {Score}\\
\midrule
Example A & 120 & 0.385 +- 0.006 & 0.759\\
Example B & 98 & 0.608 +- 0.051 & 0.206\\
\bottomrule
\end{tabularx}
\par\smallskip\footnotesize\raggedright
All values are invented. Define the actual units, uncertainty and missing
values when replacing the table.
\end{table}

\subsection{Equations, where useful}
Define all variables and assumptions before use. Equations~\eqref{eq:score}
and~\eqref{eq:difference} are layout examples only; replace or remove them.
\begin{align}
 s &= \sum_{i=1}^{n} w_i x_i,\label{eq:score}\\
 \Delta s &= s_{\mathrm{A}}-s_{\mathrm{B}}.\label{eq:difference}
\end{align}
These expressions do not represent a finding or method of the review.

\section{Critical discussion and limitations}\label{sec:discussion}
Synthesize agreements and disagreements across the referenced literature.
Describe actual uncertainties, limitations and evidence gaps, distinguishing
established observations from the authors' interpretation.

\section{Future directions}\label{sec:future}
Identify research directions that follow from the evidence and explain their
limitations. Use the authors' own judgment and avoid speculative claims
presented as established results.

\section{Conclusion}\label{sec:conclusion}
Summarize the review's central message and its implications in terms justified
by the cited work. Confirm the final sections and submission materials against
the journal's current author instructions.

\section*{Declaration of competing interests}
Replace with the authors' confirmed statement.
\section*{Acknowledgements and funding}
Replace with confirmed contributions, funding sources and grant identifiers,
or remove statements that do not apply.
\section*{Supplementary material}
Describe any supporting files that will accompany the manuscript. The optional
\texttt{supplementary.tex} compiles independently with S-numbered displays.

\bibliographystyle{elsarticle-num}
\bibliography{references}
\end{document}
